\chapter{Fundamentação teórica}
\label{cap:fundamentacao}

Este capítulo reúne os conceitos que sustentam as decisões de projeto
descritas nos capítulos seguintes: jogos sérios e aprendizagem; o domínio da
automação industrial como contexto educacional; grandes modelos de linguagem e
inferência local; a integração de IA em jogos; as ferramentas de
desenvolvimento utilizadas; e os princípios de verificação adotados.

\section{Jogos sérios e aprendizagem}
\label{sec:jogos_serios}

O termo \emph{jogos sérios} designa jogos cujo propósito principal não é o
entretenimento, e sim um objetivo explícito --- educacional, de treinamento ou
de simulação \cite{abt1970serious}. Décadas de pesquisa em aprendizagem com
jogos mostraram que essa mídia pode incorporar princípios pedagógicos de forma
orgânica: identidade do jogador, experimentação ativa, risco controlado,
problemas bem ordenados e retroalimentação imediata \cite{gee2003what}. A
transição da simulação para a realidade virtual e desta para os jogos
consolidou a ideia de que ambientes virtuais interativos podem treinar
procedimentos antes que o aprendiz os execute no mundo físico
\cite{zyda2005from}.

Revisões do campo organizam os jogos sérios por aplicação (saúde, defesa,
educação, planejamento urbano, treinamento profissional) e destacam que a
avaliação de impacto segue sendo um desafio metodológico
\cite{laamarti2014overview}. Essa observação é relevante para o presente
trabalho: a validação aqui realizada é técnica --- funcionalidade, desempenho,
latência, reprodutibilidade --- e não uma medição de aprendizagem. A separação
entre o que foi verificado e o que permanece como hipótese educacional é
deliberada e está declarada nas limitações (ver \autoref{cap:discussao}).

\section{Automação industrial como domínio educacional}
\label{sec:automacao}

A automação industrial é um domínio de ensino em que a distância entre teoria
e prática cobra um preço alto. Conceitos como lógica de intertravamento,
estados de equipamento, diagnóstico de falhas e leitura de diagramas
\emph{ladder} tornam-se concretos apenas quando o estudante os vê operando em
um sistema físico. A norma IEC~61131-3 padroniza as linguagens de programação
de controladores lógicos programáveis, incluindo a linguagem de diagrama de
contatos (\emph{ladder}) \cite{iec61131}; o trabalho educativo com CLPs gira,
em grande parte, em torno dessas linguagens e da lógica de segurança que elas
expressam.

Bancadas didáticas completas, porém, são caras, ocupam espaço, exigem
manutenção e impõem restrições de segurança. Uma simulação interativa não
substitui a bancada --- não forma habilidades motoras nem expõe o estudante
aos riscos reais ---, mas pode cumprir dois papéis complementares: preparar o
estudante para o contato com o equipamento (familiaridade com nomenclatura,
sequências de partida, leitura de sinais) e permitir prática adicional de
diagnóstico em cenários de falha que seriam inconvenientes ou perigosos de
reproduzir no laboratório físico. O Laboratório 404 foi projetado com esse
papel: seu cenário industrial ficcional pede que o jogador localize
componentes, instale módulos em um painel, reinicie um CLP, diagnostique uma
bomba e restabeleça um intertravamento de porta.

\section{Grandes modelos de linguagem}
\label{sec:llm}

Grandes modelos de linguagem são redes neurais de grande porte treinadas para
prever o próximo elemento de uma sequência de texto. A arquitetura
\emph{Transformer}, baseada em mecanismos de atenção e publicada em 2017,
viabilizou paralelismo e escala e é o substrato da maior parte dos modelos
contemporâneos \cite{vaswani2017attention}. Em uso, o modelo recebe um
\emph{prompt} (instruções e contexto) e gera uma resposta token a token; a
qualidade da resposta depende do treinamento, do tamanho do modelo e da
construção do prompt.

Duas características desses modelos importam diretamente para um jogo. A
primeira é que eles \emph{geram texto plausível}, não necessariamente verdade:
um modelo pode inventar procedimentos, componentes ou valores inexistentes
(fenômeno conhecido como alucinação). A segunda é que a geração tem custo:
quanto maior o modelo, mais memória e mais tempo de inferência. O tempo total
percebido pelo usuário decompõe-se em parcelas que este trabalho mede e
reporta separadamente:

\begin{equation}
\label{eq:latencia}
t_{\text{resposta}} = t_{\text{rede}} + t_{\text{fila}} + t_{\text{inferência}} + t_{\text{validação}} \quad [\text{s}],
\end{equation}

\noindent em que $t_{\text{rede}}$ cobre a comunicação HTTP local,
$t_{\text{fila}}$ aguarda recursos do modelo, $t_{\text{inferência}}$ é o
tempo de geração no modelo e $t_{\text{validação}}$ agrega a validação de
contrato e de intenções. Em um ambiente de jogo, se qualquer parcela crescer
além do aceitável, o sistema precisa de uma saída segura --- no caso deste
trabalho, a resposta de \emph{fallback}.

Em um contexto educacional, a alucinação deixa de ser um detalhe curioso e
passa a ser um risco de projeto. A resposta adotada aqui é de engenharia:
restringir a superfície de ação do modelo. A IA pode conversar livremente
dentro do tema, mas sua saída é reduzida a texto e a um rótulo de intenção
pertencente a um conjunto finito; nenhuma resposta do modelo é executada como
comando.

\section{Integração de IA em jogos}
\label{sec:ia_jogos}

O levantamento de Gallotta \emph{et al.} mapeia os papéis que LLMs podem
assumir em jogos --- geração de conteúdo, personagens jogáveis e não jogáveis,
apoio ao projeto, testes --- e discute limites práticos como custo, latência e
confiabilidade \cite{gallotta2024largelanguage}. Do lado da pesquisa em
agentes, \emph{Generative Agents} mostrou que agentes que registram
experiências em linguagem natural, refletem e planejam produzem comportamento
social emergente e verossímil \cite{park2023generative}; \emph{Voyager}
demonstrou um agente que aprende continuamente em um jogo aberto consultando
um modelo externo e escrevendo código \cite{wang2023voyager}. Ambos os
trabalhos exploram a \emph{autonomia}: quanto mais o modelo decide e age, mais
rico (e mais imprevisível) o comportamento resultante.

O presente trabalho investiga a hipótese complementar: que uma IA
\emph{diegética e consultiva} --- que participa da ficção, comenta o estado do
mundo e responde a perguntas, sem agir --- oferece valor narrativo e percebido
suficiente para um jogo educacional, com custo de risco muito menor. Essa
posição é compatível com a leitura do levantamento citado, que identifica nas
limitações das LLMs exatamente os pontos que um projeto educacional precisa
controlar: confiabilidade, previsibilidade e custo \cite{gallotta2024largelanguage}.

\section{Inferência local}
\label{sec:inferencia_local}

Executar o modelo na própria máquina elimina dependência de serviços externos,
dispensa chaves de aplicação e mantém o conteúdo local; o Ollama é uma
ferramenta consolidada para esse cenário, capaz de executar modelos abertos em
hardware de consumidor \cite{ollama2026}. As compensações são conhecidas:
modelos menores geram respostas em escalas de segundos, o que restringe seu
uso a interações conversacionais por turnos em vez de reações em tempo real.
Para um jogo educacional, essa restrição é compatível com o formato adotado:
o diálogo com a IA acontece em um terminal, com pausa natural para reflexão.

\section{Desenvolvimento de jogos com Godot e Blender}
\label{sec:godot_blender}

O motor Godot organiza o jogo em cenas e nós, com a linguagem GDScript para a
lógica; a documentação oficial descreve tanto o modelo de cenas quanto o
sistema de \emph{input} por ações nomeadas, que permite mapear teclado e
controles analógicos sem acoplar o código a índices de botão
\cite{godot2026}. O Blender cobre a produção dos assets 3D --- modelagem,
materiais e exportação --- e o formato glTF~2.0, em sua variante binária GLB,
foi projetado como formato de transmissão e carga de cenas 3D em tempo real
\cite{blender2026,khronos2026gltf}.

Para um jogo executável em computador pessoal sem GPU dedicada, três decisões
de projeto são determinantes: controlar a complexidade geométrica e o número
de materiais; preferir um renderizador compatível com hardware modesto (no
Godot, a estratégia \texttt{gl\_compatibility}); e medir o desempenho em
quadros por segundo na máquina-alvo. O detalhamento dessas decisões no
projeto está no \autoref{cap:desenvolvimento}.

\section{Serviços locais e contratos de dados}
\label{sec:servicos}

Manter a IA em um processo separado do jogo tem uma consequência de engenharia
direta: a comunicação entre eles precisa de um contrato explícito. O arcabouço
FastAPI permite expor um serviço HTTP local com validação automática de dados
de entrada e saída; a biblioteca Pydantic descreve e valida os modelos de
dados em Python \cite{fastapi2026,pydantic2026,python2026}. O servidor Uvicorn executa
esse serviço \cite{uvicorn2026} e a comunicação é feita por requisições HTTP
com JSON, um formato de intercâmbio textual simples e legível.

A adoção de um contrato validado nos dois lados --- adaptador e cliente no
jogo --- é o que permite tratar a IA como um componente potencialmente
falível: entradas dentro de limites, saídas dentro de um esquema conhecido e
respostas de erro previsíveis. Essa arquitetura de fronteira é o coração do
desenho de segurança adotado neste trabalho.

\section{Verificação, validação e processo}
\label{sec:verificacao}

Verificar é responder se o sistema foi construído conforme a especificação;
validar é responder se ele resolve o problema proposto. A prática de escrever
testes antes ou junto da implementação --- popularizada pelo desenvolvimento
orientado a testes \cite{beck2002tdd} --- é uma forma de manter a verificação
próxima do código. Pesquisas em engenharia de software de sistemas
informacionais propõem, para artefatos tecnológicos, o ciclo
projeto--construção--avaliação, em que cada afirmação sobre o artefato é
acompanhada de evidência \cite{hevner2004design}. Para avaliação experimental
em engenharia de software, é prática corrente explicitar hipóteses, variáveis,
ameaças à validade e limitações \cite{wohlin2012experimentation}.

Este trabalho adota uma variante pragmática desses princípios, adequada a um
desenvolvimento conduzido por uma pessoa com apoio de agentes de IA: uma
especificação com requisitos e critérios de aceitação verificáveis, um ciclo
incremental em que cada etapa termina demonstrável, uma suíte automatizada
executável por linha de comando e um registro de evidências que inclui imagens
e medições. A formalização desse processo, incluindo a matriz de
rastreabilidade, está no \autoref{cap:metodos} e no
\autoref{ap:requisitos}.

\section{Simulação, emulação e laboratórios virtuais}
\label{sec:simulacao}

Ambientes virtuais de treinamento ocupam um espectro que vai da simulação
aproximada à emulação fiel. Na \emph{simulação}, reproduz-se o comportamento
relevante de um sistema com as simplificações necessárias para tornar o modelo
computacionalmente viável; na \emph{emulação}, procura-se reproduzir o
funcionamento do equipamento original com alto grau de fidelidade, incluindo
tempos e protocolos. A escolha entre os dois polos depende do objetivo de
aprendizagem: para familiarização com procedimentos e diagnóstico, a simulação
costuma bastar; para validação de programas de CLP, a emulação é o alvo
[REFERÊNCIA NECESSÁRIA: revisão sobre emuladores de CLP para ensino].

O conceito contemporâneo de gêmeo digital --- um modelo virtual conectado a um
sistema físico, alimentado por seus dados --- amplia essa discussão, mas
também a desloca: o gêmeo digital pressupõe o sistema físico existente. Um
laboratório virtual educacional, como o proposto neste trabalho, opera em
sentido inverso: existe antes e independentemente do equipamento, para
preparar o estudante que um dia o operará [REFERÊNCIA NECESSÁRIA: estudos de
laboratórios virtuais e remotos na educação em engenharia].

Essa distinção importa para a interpretação dos resultados deste trabalho: o
Laboratório 404 é uma \emph{simulação ficcional de processos industriais}, não
um emulador de CLP. Sua fidelidade é de procedimento e de nomenclatura ---
fusível, módulo, painel, CLP, bomba, intertravamento de porta ---, e não de
protocolo ou de tempo real.

\section{Engenharia de requisitos e especificação}
\label{sec:requisitos}

A engenharia de requisitos trata da descoberta, documentação e manutenção do
que um sistema deve fazer (requisitos funcionais) e de como deve se comportar
em termos de qualidade e restrição (requisitos não funcionais)
[REFERÊNCIA NECESSÁRIA: fonte canônica de engenharia de requisitos]. Em
projetos pequenos e solitários, a tentação de dispensar esse registro é
grande; o custo aparece depois, na forma de retrabalho e de afirmações não
verificáveis sobre o produto.

Neste trabalho, adotou-se uma prática de especificação leve, porém completa
para os fins de verificação: cada requisito recebeu identificador estável; os
requisitos não funcionais foram escritos como metas mensuráveis (por exemplo,
executar em GPU integrada a uma taxa mínima de quadros; responder a uma
mensagem da IA dentro de um limite de tempo); e cada requisito funcional
ganhou um critério de aceitação observável, no formato DADO/QUANDO/ENTÃO. O
formato força a explicitação de três coisas que, de outro modo, ficariam
implícitas: a pré-condição, o estímulo e o resultado esperado --- inclusive em
situações de erro, que costumam ser as mais negligenciadas.

A matriz de rastreabilidade fecha o ciclo ao associar cada requisito ao teste
que o verifica e à evidência que o comprova. É essa matriz que permite
responder, a qualquer momento, à pergunta central de uma defesa técnica:
\emph{como você sabe que isso funciona?}

\section{Considerações do capítulo}
\label{sec:consideracoes_fundamentacao}

A fundamentação deste trabalho combina três campos que costumam caminhar
separados: a aprendizagem por jogos, a engenharia de software de sistemas
interativos e a tecnologia de modelos de linguagem. É dessa combinação que
surge o desenho investigado nos capítulos seguintes: um jogo sério 3D, com
cenário de automação industrial, cuja IA local é deliberadamente limitada a
conversar --- e é justamente essa limitação que torna o conjunto previsível o
bastante para uso educacional.
